\section{Profesión y sociedad: identidad de Builder, dolor a corto plazo y beneficio a largo plazo}

Sobre si «los programadores serán reemplazados», la respuesta de Peter no es tibia: la escasez de la programación puramente manual va a disminuir. Pero enfatiza que la capacidad de construir no desaparece, solo se reorganiza en otra capa.

\subsection{Cómo leer la frase «programar es como knitting»}

Esta metáfora fácilmente provoca rechazo, pero su contenido técnico es: \textbf{el valor marginal de la capa de implementación disminuye, mientras que el valor de la capa de definición de problemas y diseño de sistemas aumenta}. No es «ya no se necesitan personas», sino «las personas que se necesitan hacen algo distinto».

\begin{knowledgebox}{La ruta de migración de la identidad de Builder}
\begin{itemize}
    \item De «dominar la sintaxis» a «modelar el problema»;
    \item De «volumen de código» a «responsabilizarse por el resultado del sistema»;
    \item De «implementación individual» a «orquestación de la colaboración entre humano y máquina».
\end{itemize}
\end{knowledgebox}

\subsection{La doble realidad a nivel social}

El tramo final de la entrevista ofrece dos tipos de señales a la vez: por un lado, la ansiedad real sobre el empleo y la identidad; por otro, la automatización real de las pequeñas empresas y el empoderamiento real de los usuarios con discapacidad. Ambas deben sostenerse a la vez para que la discusión esté completa.

\begin{importantbox}{Una realidad que ni la política ni el producto pueden evadir}
\begin{enumerate}
    \item El desempleo estructural a corto plazo y el desajuste de competencias se agravarán;
    \item El aumento de la productividad a largo plazo creará nuevos puestos, pero con retraso en el tiempo;
    \item El sistema educativo debe pasar de «enseñar herramientas» a «enseñar pensamiento sistémico»;
    \item La responsabilidad de las empresas debe ampliarse de «mejorar la eficiencia» a «reducir el dolor de la transición».
\end{enumerate}
\end{importantbox}

\begin{figure}[H]
\centering
\includegraphics[width=0.9\textwidth]{frames/frame-025932.jpg}
\caption{Conflicto de plataformas y demandas de los usuarios: el acceso del Agent se convierte en un nuevo punto de disputa\protect\footnotemark}
\end{figure}
\footnotetext{Intervalo de tiempo en el video: 02:58:20--03:00:10.}

\subsection{Resumen de la sección}
La controversia social de OpenClaw nos recuerda que, si la narrativa tecnológica solo habla de que «será mejor», pierde credibilidad. Una discusión seria debe enfrentar a la vez las ganancias de eficiencia y el costo de la transición.

